% page 252 of the facsimile (image p-251 of the scan)
% block 162.6 x 242.0 mm ; left margin 8.0 mm
\pgc{-12.700mm}{0.000mm}{
\l{\cel{5}244.}
\l{}
\l{\sou{-ismo}.\cel{2}Sufixo qua signifikas \textquotedbl{}doktrino\textquotedbl{}, \textquotedbl{}partiso\textquotedbl{}, \textquotedbl{}sistemo\textquotedbl{}.}
\l{}
\l{\sou{-isto}.\cel{2}Sufixo qua signifikas : I.\textquotedbl{}la persono qua profesione}
\l{\cel{3}su okupas pri...\textquotedbl{}. II. adepto, adherinto di partiso, di}
\l{\cel{3}skolo doktrinala.}
\l{}
\l{\sou{istmo}. (\sou{geogr.}) Streta lango de tero, inter du mari, e qua un-}
\l{\cel{3}ionas peninsulo a la kontinento cetera. - DEFIS.}
\l{}
\l{\sou{-it.}\cel{3}Dezinenco di la participo pasintala pasiva.}
\l{}
\l{\sou{ita}.\cel{2}La persono, la kozo, qua maxime distas de ta qua parolas;}
\l{\cel{3}la persono, la kozo, quan onu nomis enumer-unesme.}
\l{}
\l{\sou{+itemo.}\cel{2}Singla ek la frazi (generale kun nombri) en enumero,}
\l{\cel{3}etato, budjeto, fakturo, kelkulo. - EFI.}
\l{}
\l{\sou{iterar}. (\sou{trans}.) Ri-agar lo sama. - EFIS.}
\l{}
\l{\sou{itinerario}.\cel{2}I. Indiko pri la voyo uzenda.-.II. Indiko pri omna}
\l{\cel{3}loki ube onu pasas por irar de lando ad altra, qua kontenas}
\l{\cel{3}ulakaze deskripti di la loki, o noti e voyajo-memori. -}
\l{\cel{3}III. Indiko di omna stacioni qui existas sur la parkuro di}
\l{\cel{3}relvoyo, di la diversa treni qui pasas en oli ; di la trans-}
\l{\cel{3}porto-preco, e c. - DEFIS.}
\l{}
\l{\sou{-iv.}\cel{2}Sufixo qua signifikas (soldite a radiko verbala) \textquotedbl{}kapabla,}
\l{\cel{3}povanta, apta por...\textquotedbl{}}
\l{}
\l{\sou{-ivoro}.\cel{2}Sufixo qua signifikas \textquotedbl{}qua su nutras ye...\textquotedbl{}.}
\l{}
\l{\sou{ivoro}. I, Materio ye qua konsistas la dentegi di elefanto.}
\l{\cel{3}- II. Materio ye qua konsistas la denti di ula pakidermi.-}
\l{\cel{3}- EFI.}
\l{}
\l{\sou{-iz}.\cel{2}Sufixo qua signifikas \textquotedbl{}provizar, garnisar,indutar, im-}
\l{\cel{3}pregnar ye...\textquotedbl{}}
\l{}
\l{\sou{izabela}.\cel{2}Pale-flava. - DEFIS.}
\l{}
\l{\sou{izobaro}.\cel{2}(\sou{fiziko)}.\cel{2}Lineo ek la punti di la Terglobo ube la}
\l{\cel{3}preso atmosferala esas sama ye instanto determinita. -}
\l{\cel{3}DEFIRS.}
\l{}
\l{\sou{izocela}. (\sou{geom.})\cel{2}Di qua la lateri esas inter-egala. - EFIS.}
\l{}
\l{\sou{izodinama}. (\sou{fiziko}). Di qua la forco esas inter-egala surla du}
\l{\cel{3}lateri. - DEFIS.}
\l{}
\l{\sou{izogama}. (\sou{biol.})\cel{2}Dicesas pri la vejetanti inferiora, che qui}
\l{\cel{3}la genito-elementi, qui interunionas por formacar la ovo,}
\l{\cel{3}esas intersimila. - DEFIS.}
\l{}
\l{izogona. (\sou{geom.)}\cel{2}Di qua la anguli esas inter-egala. DEFIS.}
}
